\documentclass[10pt,a4paper]{amsart}
\usepackage[margin=19mm]{geometry}
\usepackage{amsmath,amssymb,mathtools}
\usepackage[scr=rsfs,cal=euler]{mathalfa}
\usepackage{xfrac}
\usepackage[unicode,hidelinks,bookmarks=false]{hyperref}
\newcommand{\ie}{i.e.\ }
\newcommand{\cf}{cf.\ }
\newcommand{\resp}{respectively\ }
\newcommand{\Subsection}{\subsection}
\providecommand{\nobreakdash}{\nobreak\hbox{-}}
\setlength{\emergencystretch}{2em}
\multlinegap0pt
\usepackage{bm}
\usepackage{amssymb}
\newtheorem{proposition}{Proposition}
\title{Identically vanishing $k$-generalized Fibonacci polynomials}
\author{S. R. Mane}
\date{Source: arXiv:2507.11596v4 (CC BY 4.0)}
\begin{document}
\maketitle

\noindent\emph{Source and licence.} Identically vanishing $k$-generalized Fibonacci polynomials. Version arXiv:2507.11596v4 (CC BY 4.0), distributed by arXiv under the Creative Commons Attribution 4.0 International licence, http://creativecommons.org/licenses/by/4.0/, as recorded in the arXiv licence metadata for this exact version and shown on the abstract page https://arxiv.org/abs/2507.11596v4. Full licence notice: \texttt{licenses/arxiv-2507.11596-CC-BY-4.0.txt}.\par\smallskip

\noindent\textit{Editorial scope.} Counted unit: the article's proposition prop:sec\_highest\_neg (source0.tex lines 755--765). The article's complete original proof of it is delivered with it (source0.tex lines 766--821, one \texttt{proof} environment of 56 lines). No mathematics is written, repaired or re-derived: the statement and the proof are the article's own lines, in the article's own notation, and no retained line is substituted or re-flowed.  Retained uncounted as the source-local context the proof needs: lines 388--394; lines 683--688.  Nothing was omitted from any retained span, so the package carries no omission record.  No retained line contains a citation, so the wrapper carries no bibliography and no bibliography entry.  No retained line contains a figure environment, an image-inclusion command or drawing code, so no image is delivered and the package declares no asset dependency.  Editorial formatting only, in addition: an AMS \texttt{amsart} preamble that restates the article's own declarations and macros used by the retained lines, namely the environments \texttt{proposition} on separate counters, together with the standard packages those lines need.  The retained environments print in this document as Proposition 1 in order of appearance, whereas the article prints the same items with the numbering of its own full document.  The complete native source of the article as distributed in the arXiv e-print for this exact version is bundled unchanged as \texttt{sources/181638/source0.tex}.
\begin{equation}
\label{eq:Cnk_neg_qr}    
\begin{split}
  C_k(-(1+h),|n|+1-k(1+h)) &= \sum_{s=0}^{q-h-1} (-1)^{k(q-s-h-1)+r} 
  \binom{1+h}{k(q-h-s-1)+r} \binom{h+s}{s} \,.
\end{split}
\end{equation}

\begin{equation}
\label{eq:high_term_neg}    
H_{n,k}(x) = \begin{cases} x^{|n|+1-k} &\qquad (r_{n,k}=0)\,,\phantom{\biggl|}
  \\ \displaystyle
  (-1)^{r_{n,k}} \binom{q_{n,k}-1}{r_{n,k}-1}x^{|n|+1-kr_{n,k}} &\qquad (r_{n,k}\in[1,k-1])\,. \end{cases} 
\end{equation}

\begin{proposition}\label{prop:sec_highest_neg}
For $k\ge2$ and $n\le0$, the second highest term, say $S_{n,k}(x)$, in a nonvanishing polynomial which is also not a monomial, is as follows:
\begin{equation}
\label{eq:sec_term_neg}    
S_{n,k}(x) = \begin{cases}
(2q_{n,k}-3)x^{|n|+1-2k} &\quad (r_{n,k}=0,\, k=2)\,,\phantom{\biggl|} \\
(q_{n,k}-1)x^{|n|+1-2k} &\quad (r_{n,k}=0,\, k\ge3)\,,\phantom{\biggl|} 
  \\ \displaystyle (-1)^{r_{n,k}} (r_{n,k}+1)\binom{q_{n,k}-1}{r_{n,k}}x^{|n|+1-k(r_{n,k}+1)} &\quad (r_{n,k}\in[1,k-1])\,. \end{cases} 
\end{equation}
To avoid a monomial, we require $|n|\ge 2k-1$ for $r_{n,k}=0$ and $|n|\ge k(r+1)-1$ for $r_{n,k}\in[1,k-1]$.
\end{proposition}

\begin{proof}
This is a followup of eq.~\eqref{eq:high_term_neg} and the derivation follows the same reasoning.
We again drop the subscripts on $q$ and $r$.
First suppose $r=0$.
For the highest power, we set $h=0$, hence for the second highest power we set $h=1$.
Then eq.~\eqref{eq:Cnk_neg_qr} yields the following:
\begin{equation}
\label{eq:Ck_low0}
\begin{split}
  C_k(-2,|n|+1-2k) &= \sum_{s=0}^{q-2} (-1)^{k(q-s-2)} \binom{2}{k(q-s-2)} \binom{s+1}{s} \,.
\end{split}
\end{equation}
A solution exists only if $k(q-s-2)\le2$, else the first binomial coefficient in eq.~\eqref{eq:Ck_low0} vanishes.
If $k\ge3$, we must have $s=q-2$, whence $k(q-s-2)=0$.
Then only the last term in the sum in eq.~\eqref{eq:Ck_low0} contributes and we obtain %the following:
\begin{equation}
C_k(-2,|n|+1-2k) = \binom{2}{0} \binom{q-1}{q-2} = q-1 \,.
\end{equation}
If $k=2$, there is also a solution where $s\ge q-3$, because $k(q-s-2) = 2$ if $s=q-3$ and $k(q-s-2) = 0$ if $s=q-2$. 
Then eq.~\eqref{eq:Ck_low0} yields the following:
\begin{equation}
\begin{split}
  C_k(-2,|n|+1-2k) &= \sum_{s=q-3}^{q-2} (-1)^{2(q-s-2)} \binom{2}{2(q-s-2)} \binom{s+1}{s}
  \\
  &= \binom{2}{2} \binom{q-2}{q-3} + \binom{2}{0} \binom{q-1}{q-2}
  \\
  &= 2q-3 \,.
\end{split}
\end{equation}
For all $k\ge2$, the corresponding power of $x$ is $x^{|n|+1-2k}$.
This proves the first two cases in eq.~\eqref{eq:sec_term_neg}.
Next suppose $1 \le r \le k-1$. 
Employing the same reasoning used for eq.~\eqref{eq:high_term_neg} for $r>0$,
now we deduce the lowest value of $h$ for which a solution exists is $h=r$. We obtain the following:
\begin{equation}
\label{eq:Ck_low1}
\begin{split}
  C_k(-(1+r),|n|+1-k(1+r)) &= \sum_{s=0}^{q-r-1} (-1)^{k(q-s-r-1)+r} %\times
  %\\
  %&\qquad\quad
  \binom{r+1}{k(q-r-s-1)+r} \binom{r+s}{s} \,.
\end{split}
\end{equation}
If $q<r+1$ there is no second highest term, hence we assume $q\ge r+1$.
We must have $q-r-s-1=0$, i.e.~$s=q-r-1$, else the first binomial coefficient in eq.~\eqref{eq:Ck_low1} vanishes.
Then eq.~\eqref{eq:Ck_low1} has a solution for all $k\ge2$.
Only the last term in the sum in eq.~\eqref{eq:Ck_low1} contributes and we obtain the following:
\begin{equation}
\begin{split}
  C_k(-(1+r),|n|+1-k(1+r)) &= (-1)^r \binom{r+1}{r} \binom{q-1}{q-r-1}
  \\
  &= (-1)^r (r+1) \binom{q-1}{r} \,.
\end{split}
\end{equation}
The corresponding power of $x$ is $x^{|n|+1-k(r+1)}$.
\end{proof}

\end{document}
